% 2022 新高考Ⅱ卷 3：举架结构抽象出的屋顶截面折线图
% 左半支自原点 O 经 D、C、B 升至顶点 A，右半支关于过 A 的竖直线对称（不标字母）。
% 步长恒为 1，举依次取 0.5、0.7、0.8、0.9，即 O(0,0)、D(1,0.5)、C(2,1.2)、B(3,2.0)、A(4,2.9)。
% 虚线小阶梯给出各段的举与步：垂足 D_1(1,0)、C_1(2,0.5)、B_1(3,1.2)、A_1(4,2.0)，
% 水平段 OD_1、DC_1、CB_1、BA_1 四段等长（即「步」），竖直段 DD_1、CC_1、BB_1、AA_1 即「举」。
\begin{tikzpicture}[scale=1.12, font=\small]
  % ---- 坐标轴 ----
  \draw[-stealth] (-0.45,0) -- (8.95,0);
  \draw[-stealth] (0,-0.4) -- (0,3.45);
  % ---- 举步阶梯（虚线）----
  \draw[dashed] (0,0) -- (1,0) -- (1,0.5) -- (2,0.5) -- (2,1.2)
                     -- (3,1.2) -- (3,2.0) -- (4,2.0) -- (4,2.9);
  % ---- 屋顶截面折线 ----
  \draw[line width=0.9pt]
    (0,0) -- (1,0.5) -- (2,1.2) -- (3,2.0) -- (4,2.9)
          -- (5,2.0) -- (6,1.2) -- (7,0.5) -- (8,0);
  % ---- 点 ----
  \foreach \p in {(0,0),(1,0.5),(2,1.2),(3,2.0),(4,2.9),(1,0),(2,0.5),(3,1.2),(4,2.0)}
    { \fill \p circle (1.6pt); }
  % ---- 标注：折线点在左上，垂足在下/右 ----
  \node[below left,  inner sep=2pt] at (0,0)   {$O$};
  \node[above left,  inner sep=1pt] at (1,0.5) {$D$};
  \node[above left,  inner sep=1pt] at (2,1.2) {$C$};
  \node[above left,  inner sep=1pt] at (3,2.0) {$B$};
  \node[above,       inner sep=2pt] at (4,2.9) {$A$};
  \node[below,       inner sep=2pt] at (1,0)   {$D_1$};
  \node[below right, inner sep=1pt] at (2,0.5) {$C_1$};
  \node[below right, inner sep=1pt] at (3,1.2) {$B_1$};
  \node[below right, inner sep=1pt] at (4,2.0) {$A_1$};
\end{tikzpicture}
